واصل الدولار الأمريكي اتجاهه الهبوطي، متراجعا أمام العملات العالمية الرئيسية في ظل توقعات تراجع معدلات التضخم وتباين توجهات السياسات النقدية. وقد أسهم هذا التراجع في تيسير الأوضاع المالية في الاقتصادات المتقدمة، وخفف من قيود التمويل الخارجي ودعم تدفقات رأس الأموال إلى الأسواق الناشئة.\textsuperscript{5}

استمرت أسواق السلع في إظهار اتجاهات متباينة، نتيجة تغير أوضاع العرض والطلب وتقلب المعنويات السوقية عبر مجموعات السلع الرئيسية، حيث شهدت أسعار الطاقة تراجعا ملحوظا، في حين حققت المعادن الثمينة وبعض السلع الزراعية مكاسب كبيرة. وانخفضت أسعار الغاز الطبيعي وسط توقعات ضعف الطلب ووفرة المعروض. وفي المقابل، ارتفعت أسعار الذهب مدفوعة بتزايد حالة عدم اليقين، واستمرار توجه المستثمرين نحو الأصول الملاذ الآمن. أما الأسواق الزراعية فجاء أداؤها متباينا، حيث ارتفعت أسعار الأرز واللحوم نتيجة تزايد ضغوط العرض.\textsuperscript{6}
